5.6 Examples of Differential Calculus in Natural Science 295

It can be seen from (5.151) that for $\alpha > 0$ a body falling in the atmosphere reaches
a steady state at which $v(t) \approx \frac{m}{\alpha}g$. Thus, in contrast to fall in airless space, the
velocity of descent in the atmosphere depends not only on the shape of the body, but
also on its mass. As $\alpha \to 0$, the right-hand side of (5.151) tends to $v_0 + gt$, that is,
to the solution of Eq. (5.149) obtained from (5.150) when $\alpha = 0$.
Using formula (5.151), one can get an idea of how quickly the limiting velocity
of fall in the atmosphere is reached.
For example, if a parachute is designed to that a person of average size will fall
with a velocity of the order $10$ meters per second when the parachute is open, then,
if the parachute opens after a free fall during which a velocity of approximately
$50$ meters per second has been attained, the person will have a velocity of about
$12$ meters per second three seconds after the parachute opens.
Indeed, from the data just given and relation (5.151) we find $\frac{mg}{\alpha} \approx 10$, $\frac{m}{\alpha} \approx 1$,
$v_0 = 50$ m/s, so that relation (5.151) assumes the form
$$v(t) = 10 + 40e^{-t}.$$
Since $e^3 \approx 20$, for $t=3$, we obtain $v \approx 12$ m/s.
Note that at the opening of the parachute, which is significantly bigger on the
front surface compared with the body of the skydiver, the force of resistance is
already proportional to the square of the speed. This means that after the opening of
the parachute, the slowdown will be quicker than that obtained from the calculation
with formula (5.151), corresponding to Eq. (5.150).
It is useful to rewrite and solve the equation of type (5.150) with this new assump-
tion of the quadratic dependence of the resistance force on the speed of movement
and to make the necessary correction to the result obtained in the first calculation.$^{32}$

\subsection{5.6.5 The Number $e$ and the Function $\exp x$ Revisited}
By examples we have verified (see also Problems 3 and 4 at the end of this section)
that a number of natural phenomena can be described from the mathematical point
of view by the same differential equation, namely
$$f'(x) = \alpha f(x), \quad (5.152)$$
$^{32}\text{It should be noted that here, like in many others studied examples of problems of natural science,}$
$\text{we illustrate the calculation method. In the calculation are included some original data (empirical}$
$\text{constants, various assumptions) on which depends the level of concordance of the calculations with}$
$\text{the reality. This data are not given by mathematics but by a corresponding science. For example,}$
$\text{in order to obtain data on the aerodynamic properties of the bodies, in particular the nature of the}$
$\text{change in the resistance force occurring during their movement in the atmosphere, it is necessary}$
$\text{often to perform a long series of experiments in wind tunnels. Mathematics helps to calculate one}$
$\text{or another mathematical model of the phenomenon, it often provides a ready calculation procedure,}$
$\text{but the correctness of the initial data and the adequacy of the model is the work of another relevant}$
$\text{science (biology, chemistry, physics, ...).}$